\section{Post-training}\label{sec:post}
\subsection{Thinker}

The post-training phase comprises a three-stage training process for Thinker, enabling \method to possess instruction-following capabilities. The dataset, designed in the ChatML~\citep{chatml} format, includes pure text-based dialogue data, visual modality conversation data, audio modality conversation data, and mixed-modality conversation data.

In the first stage, we introduce a lightweight Supervised Fine-Tuning (SFT) to bridge the gap between pretrained representations and downstream task requirements through targeted instruction optimization. SFT deliberately diverges from the pretraining data schema while maintaining architectural consistency with the pretrained model, enabling efficient knowledge transfer and preserving the completeness of the pretrained features.

The second stage adopts the Strong-to-Weak Distillation pipeline as described in Qwen3~\citep{qwen3} to further improve model performance. This distillation process consists of two main phases:

\begin{enumerate}[label=(\arabic*)]
    \item \textbf{Off-policy Distillation}: In the initial phase, outputs generated by teacher models are combined to provide response distillation. This helps lightweight student models acquire fundamental reasoning abilities, establishing a strong foundation for subsequent on-policy training.
    \item \textbf{On-policy Distillation}: In the second phase, the student model generates the responses based on sampled prompts. These on-policy sequences are then used for fine-tuning, where the student’s predicted logits are aligned with those of a teacher model (Qwen3-32B or Qwen3-235B-A22B) by minimizing the KL divergence.
\end{enumerate}





Finally, we leverage GSPO~\citep{gspo} to comprehensively enhance the model's capabilities and stability across various modalities, including text, image, video, and audio. To provide feedback for the aforementioned modalities, we employ two different types of rewards:

\begin{itemize}
    \item \textbf{Rule-based Reward}: For verifiable multimodal tasks (e.g., mathematics, coding, instruction following), the reward signal is derived from a set of predefined rules. Well-designed rule-based rewards can assess the correctness of model outputs with high precision, preventing issues like reward hacking.

    \item \textbf{Model-based Reward}: To assess performance on multimodal tasks that lack objective, predefined evaluation metrics, we adopt an LLM-as-a-judge protocol. The role of the automated evaluator is filled by Qwen3 for general tasks, while the specialized vision-language model, Qwen2.5-VL, is used for visually-grounded tasks. To ensure a more robust and grounded assessment, the LLM evaluator is furnished with the corresponding ground-truth or reference answer for a given query, where applicable.
\end{itemize}





% For multimodal tasks where objective evaluation via predefined rules is not feasible, we employ a large language model (LLM), such as Qwen3, to assess the quality of our model's responses. For tasks with a visual component, this evaluation is specifically conducted by Qwen2.5-VL. In instances where a ground-truth or reference answer exists for a query, it is also supplied to the evaluator for reference.

\subsection{Talker}
We introduce a four-stage training process for Talker, enabling Qwen3-Omni to generate speech response simultaneously with text. All training data is structured in the ChatML format to ensure consistency with Thinker.

In the first stage, we leverage hundreds of millions of speech data with multimodal context to train Talker, establishing a monotonic mapping from multimodal representation to speech. In the second stage, we perform Continual Pretraining (CPT) with high-quality data, which alleviates hallucinations caused by noisy data in the first stage and significantly improve the quality of generated speech. Concurrently, we perform long-context training that enhances Talker’s ability to process extended and complex inputs and generate contextually appropriate speech response. In the third stage, to improve the generalization of multilingual speech generation and system stability, we construct preference pairs from diverse multilingual speech samples and optimize the model using Direct Preference Optimization (DPO)~\citep{rafailov2024direct}. Finally, we apply speaker fine-tuning on the aforementioned base model, enabling Talker to adopt specific voices while refining the naturalness, expressiveness, and controllability of its speech response.

% \jin{We introduced a three-stage training process for Talker, allowing \method to generate text and speech responses simultaneously. In the first stage, we train Talker to learn context continuation. The second stage utilized DPO~\citep{rafailov2024direct} to enhance the stability of speech generation. In the third stage, we applied multi-speaker instruction fine-tuning to improve the naturalness and controllability of the speech responses.

% During the In-Context Learning (ICL) training phase, in addition to utilizing text supervision similar to that of Thinker, we perform a speech continuation task through next-token prediction, leveraging an extensive dataset of dialogues that incorporate multimodal contexts and spoken responses. Talker learns to establish a monotonic mapping from semantic representation to speech, while also acquiring the ability to express speech with diverse attributes that are contextually appropriate, such as prosody, emotion, and accent. Additionally, we implement timbre disentanglement techniques to prevent the model from associating specific voices with infrequent textual patterns.

% \begin{equation}
%     \mathcal{L}_{\text{DPO}}(\mathcal{P}_\theta; \mathcal{P}_{\text{ref}}) = -\mathbb{E}_{(\bm{x}, \bm{y_w}, \bm{y_l}) \sim \mathcal{D}} \left[ \log \sigma \left( \beta \log \frac{\mathcal{P}_\theta(\bm{y_w} \mid \bm{x})}{\mathcal{P}_{\text{ref}}(\bm{y_w} \mid \bm{x})} - \beta \log \frac{\mathcal{P}_\theta(\bm{y_l} \mid \bm{x})}{\mathcal{P}_{\text{ref}}(\bm{y_l} \mid \bm{x})} \right) \right].
% \end{equation}

% To broaden the coverage of speakers and scenarios, the pretraining data inevitably contains label noise and pronunciation errors, leading to model hallucinations. 
% To mitigate this issue, we introduce a reinforcement learning phase to improve the stability of speech generation.  Specifically, for each request and response text paired with the reference speech, we build a dataset $\mathcal{D}$ with the triplet data $(\bm{x}, \bm{y_w}, \bm{y_l})$, where $\bm{x}$ is the input sequence with input text, and $\bm{y_w}$ and $\bm{y_l}$ are the good and bad generated speech sequences respectively. We rank these samples based on their reward scores associated with word error rate~(WER) and the punctuation pause error rate.

% Lastly, we performed speaker fine-tuning on the aforementioned base model, enabling Talker to adopt specific voices and improve its naturalness.} %produce style-controllable responses. 
%We employed highly expressive data from multiple speakers, assigning a unique index token to each speaker and substituting the <|talker\_bos|> tokens accordingly. The experimental results showed that the speaker fine-tuned model could more precisely capture the nuanced prosodic styles of target speakers while preserving the foundational stability provided by the base model.
%To further improve Talker's controllability, we integrated additional instructed dialogue data, enabling \method to comprehend user intentions and flexibly manage various aspects of speech responses, such as rate, style, emotion, dialect, and role-playing.





\subsection{Captioner}
Captioning is a foundational task in multimodal understanding, integral to the training and evaluation of large multimodal models. However, the vast majority of existing research has concentrated on visual captioning, largely neglecting the audio modality. This omission is significant, as auditory perception is a crucial component of human sensory experience and interaction with the world. To address this gap and facilitate more comprehensive research in multimodal perception, we introduce the Qwen3-Omni-30B-A3B-Captioner. This model was developed by fine-tuning the Qwen3-Omni-30B-A3B on a large-scale dataset of detailed audio descriptions. The resulting system generates detailed, low-hallucination captions for arbitrary audio inputs. The \textbf{Appendix}~\ref{sec:caption-case} provides qualitative results that demonstrate our model's captioning capabilities across diverse acoustic scenarios.




